% Preamble for CMF research papers.
% Defines the document class, fonts, theorem styles, and unicode mappings.
% Used by paper1.tex / paper2.tex / paper3.tex (which inline this content).
% To regenerate a self-contained paperN.tex from a body file:
%   cat preamble.tex paperN.body.tex > paperN.tex
% Compile with: xelatex paperN.tex
%
% This preamble does NOT include the \title/\author commands —
% those are at the top of each paperN.body.tex (paper-specific).

\documentclass[11pt]{article}

\usepackage[margin=1in]{geometry}
\usepackage{amsmath,amssymb,amsthm}
\usepackage{mathtools}
\usepackage{bm}

% Use xelatex with TeX Gyre Termes (Times equivalent) for body and math
\usepackage{fontspec}
\usepackage{unicode-math}
\setmainfont{TeX Gyre Termes}
\setmonofont{DejaVu Sans Mono}[Scale=0.85]
\setmathfont{Latin Modern Math}[Scale=MatchLowercase]

\usepackage{textcomp}
\usepackage{booktabs}
\usepackage{array}
\usepackage{longtable}
\usepackage{calc}
\usepackage{etoolbox}
\usepackage{makecell}
\usepackage[hidelinks]{hyperref}
\usepackage{enumitem}
\usepackage{xcolor}
\usepackage{framed}
\usepackage{fvextra}

% Enable pandoc's syntax-highlighted code-block environment.

% Tighter spacing
\setlength{\parskip}{0.5\baselineskip}
\setlength{\parindent}{0pt}

% Section heading style: bold, smaller than default. The markdown text already
% includes the section number ("## 1. Introduction"), so we suppress LaTeX's
% auto-numbering here.
\usepackage{titlesec}
\titleformat{\section}[block]{\normalsize\bfseries}{}{0pt}{}
\titleformat{\subsection}[block]{\normalsize\bfseries}{}{0pt}{}
\titleformat{\subsubsection}[block]{\normalsize\bfseries}{}{0pt}{}
\titlespacing*{\section}{0pt}{1.2\baselineskip}{0.4\baselineskip}
\titlespacing*{\subsection}{0pt}{1.0\baselineskip}{0.3\baselineskip}

% Pandoc-required commands
\providecommand{\tightlist}{\setlength{\itemsep}{0pt}\setlength{\parskip}{0pt}}
\providecommand{\pandocbounded}[1]{#1}

% Custom title block (centered, matching v2 style)
\makeatletter
\renewcommand\maketitle{%
  \begingroup
  \begin{center}
    \LARGE\bfseries \@title \par
  \end{center}
  \ifx\@subtitle\@empty\else
    \vspace{0.4em}
    \begin{center}
      \itshape \@subtitle \par
    \end{center}
  \fi
  \vspace{0.6em}
  \begin{center}
    \@author \par
    \vspace{0.2em}
    Correspondence: \@correspondence \par
    \vspace{0.2em}
    \@date \par
  \end{center}
  \vspace{1em}
  \endgroup
}
\def\@subtitle{}
\newcommand{\subtitleX}[1]{\def\@subtitle{#1}}
\def\@correspondence{}
\newcommand{\correspondenceX}[1]{\def\@correspondence{#1}}
\makeatother

% Custom abstract: bold heading centered, indented body
\renewenvironment{abstract}{%
  \begin{center}\bfseries Abstract\end{center}%
  \begin{quote}\setlength{\parindent}{0pt}}{\end{quote}\medskip}

% Keywords command
\newcommand{\keywordsX}[1]{%
  \begin{quote}\setlength{\parindent}{0pt}\textbf{Keywords:} #1\end{quote}}

\begin{document}
% Paper-specific metadata
\title{On the Tangent of the Argument: Closed Forms for Constants
Arising from Anticommutative Pauli CMF Matrix Products}

\author{Louis Blaine}
\correspondenceX{louisblaine45@gmail.com}
\date{March 2026}

% Body
\maketitle

\begin{abstract}
In a companion paper {[}1{]} we introduced anticommutative Conservative
Matrix Fields (CMFs) on the Gaussian integer lattice whose step matrices
lie in the Pauli algebra, and identified five numerical constants
produced by matrix products that could not be expressed in closed form.
Here we provide exact closed forms for all five. The key result is a
theorem connecting the Pauli CMF recurrence to an alternating arctangent
series via a Möbius angle recursion: the constant equals
\textbar tan(S)\textbar{} where S = \(\Sigma\)(−1)\^{}(j+1)
arctan(p(j)/r(j)). All five constants share the unified form c =
\textbar tan(arg(\(\Gamma\)(1+w/2)/\(\Gamma\)(1/2+w/2)))\textbar{} =
\textbar Im(z)/Re(z)\textbar{} for a complex parameter w and z =
\(\Gamma\)(1+w/2)/\(\Gamma\)(1/2+w/2). We further prove the exact
identity S\textsubscript{1} + S\textsubscript{3} = \(\pi\)/4, which
relates two of the five constants. We also establish that
S\textsubscript{1} is the phase of a zeta-regularised functional
determinant on the circle, and report extensive negative PSLQ searches
at the stated precision and coefficient bounds, ruling out relations to
Dirichlet L-values, polylogarithms, Bloch-Wigner functions, and Mahler
measures.
\end{abstract}

\keywordsX{Conservative Matrix Field, Pauli algebra, complex Gamma
function, Möbius transformation, arctan series, zeta-regularised
determinant, Hurwitz zeta function}

\hypertarget{introduction}{%
\section{1. Introduction}\label{introduction}}

The companion paper {[}1{]} extended the CMF framework of Raz et
al.~{[}2{]} to the Gaussian integer lattice \(\mathbb{Z}\){[}i{]},
showing that anticommutative paths admit a natural realisation in the
Pauli algebra. Five constants generated by explicit polynomial step
matrices resisted identification. The present paper provides exact
closed forms for all five, proves an exact relation between two of them,
interprets them spectrally, and reports a systematic numerical
investigation of their arithmetic nature. Section 2 establishes the
Möbius Angle Theorem. Section 3 derives exact Gamma closed forms.
Section 4 proves S\textsubscript{1}+S\textsubscript{3} = \(\pi\)/4.
Section 5 gives the spectral interpretation. Section 6 reports extensive
negative PSLQ searches.

\hypertarget{the-muxf6bius-angle-theorem}{%
\section{2. The Möbius Angle
Theorem}\label{the-muxf6bius-angle-theorem}}

\hypertarget{role-of-the-pauli-structure}{%
\subsection{2.1. Role of the Pauli
structure}\label{role-of-the-pauli-structure}}

The Pauli anticommutation \{\(\sigma _{i}\), \(\sigma _{j}\)\} =
2\(\delta _{ij}\)I constrains the CMF step matrices to trace-zero real
combinations of two Pauli generators, which in suitable coordinates take
the form M = r\(\sigma\)\_z + p\(\sigma\)\_x = {[}{[}r, p{]}, {[}p,
−r{]}{]}. Two properties of this form are consequences of the
anticommutation, not free parameters:\textsuperscript{1}

\textbf{(i) Involution.} M\textsuperscript{2} = (r\textsuperscript{2} +
p\textsuperscript{2})I, so the induced projective action z \(\mapsto\)
(rz + p)/(pz − r) satisfies f \(\circ\) f = id. This is what produces
the alternating sign in S = \(\Sigma _{j}\) (−1)\^{}(j+1)
arctan(p(j)/r(j)) below; generic Möbius iteration does not telescope
into sums of arctangents.

\textbf{(ii) Tangent linearisation.} The antidiagonal-plus-diagonal form
yields a one-step arctan increment arctan(p/r) when the argument of z is
tracked, collapsing two-dimensional matrix dynamics on
\(\mathbb{C}\mathbb{P}^{1}\) into one-dimensional additive angle
dynamics.

The anticommutation thus supplies exactly the algebraic data required
for (i) and (ii). Without it, the step matrices would produce generic
Möbius iterates whose orbits do not admit closed forms of the type
derived here. However, once the arctan sum is obtained, its reduction to
a Gamma-function ratio (Lemma 2.2 below) proceeds by classical
special-function identities and does not rely further on the Pauli
framework. The companion paper {[}4{]} correspondingly treats the
resulting Gamma-ratio formula as an independent object, with the Pauli
origin appearing only as motivation.

\textbf{Theorem 2.1.} The accumulated Pauli CMF matrix ratio z\_N
satisfies the Möbius recurrence z\_k = (r(k)z + p(k))/(p(k)z − r(k)),
conjugate to the angle recursion \(\varphi\)\emph{k = arctan(p(k)/r(k))
− \(\varphi\)}\{k−1\} via z\_k = −tan(\(\varphi\)\emph{k). The CMF
constant is c = \textbar tan(S)\textbar{} where S =
\(\Sigma\)}\{j\(\geq\)1\} (−1)\^{}(j+1) arctan(p(j)/r(j)).

\emph{The proof follows by substitution and the tangent subtraction
identity.} \(\square\)

\textbf{Lemma 2.2} (Gamma Product Formula). For Re(w) \textless{} 1:

\[\prod_{k=1}^{N} \left(1+\frac{w}{k}\right)^{(-1)^{k+1}} \to \sqrt{\pi} \cdot \frac{\Gamma(1+w/2)}{\Gamma(1/2+w/2)} \quad \text{[even-}N\text{ limit]}\]

The factor \(\sqrt{\pi}\) arises from the Weierstrass product
normalisation at the half-integer shift. Since \(\sqrt{\pi}\) is real
and positive, arg(\(\prod\)) =
arg(\(\Gamma\)(1+w/2)/\(\Gamma\)(1/2+w/2)); all downstream uses in
Section 3 depend only on the argument and are unaffected.

\hypertarget{convergence-and-precision}{%
\subsection{2.3. Convergence and
precision}\label{convergence-and-precision}}

The alternating arctangent sum S = \(\Sigma _{j}\) (−1)\^{}(j+1)
arctan(p(j)/r(j)) converges conditionally by the Leibniz criterion
(arctan(p(k)/r(k)) \(\to\) 0 monotonically in all five cases). The
truncation error after N terms satisfies \textbar S − S\_N\textbar{}
\(\leq\) arctan(p(N+1)/r(N+1)) \(\leq\) p(N+1)/r(N+1).

For c\textsubscript{1} (p=1, r=k) this bound is 1/(N+1), so direct
summation requires O(10\^{}120) terms for 120-digit precision ---
plainly impractical. In practice, all five constants are evaluated via
the closed-form Gamma ratio (Lemma 2.2 and Section 3), using mpmath's
internal Spouge approximation for \(\Gamma\) at complex arguments.
Spouge's method requires O(120 ln 10) \(\approx\) 277 summation terms
for 120-digit relative error, making high-precision evaluation
straightforward. The Weierstrass product of Lemma 2.2 converges
absolutely since \(\Sigma\)\textbar w/k\textbar{}\textsuperscript{2} =
\textbar w\textbar{}\textsuperscript{2}\(\pi^{2}\)/6 \textless{}
\(\infty\); the product-to-Gamma identification is a standard
consequence of the Weierstrass factorisation {[}7, §12.1{]}.

\hypertarget{branch-conventions}{%
\subsection{2.4. Branch conventions}\label{branch-conventions}}

We use the principal branch Arg : \(\mathbb{C}\)* \(\to\) (−\(\pi\),
\(\pi\){]} throughout. The constant c = \textbar tan(Arg(z))\textbar{}
where z = \(\Gamma\)(1+w/2)/\(\Gamma\)(1/2+w/2) admits the manifestly
branch-free equivalent form c = \textbar Im(z)/Re(z)\textbar{} (valid
when Re(z) \(\neq\) 0). Numerical evaluation at 120-digit precision
confirms Re(z) \textgreater{} 0 for all five constants (see companion
paper {[}4, §2.3{]} for the general statement). For these constants
Arg(z) \(\in\) (−\(\pi\)/2, \(\pi\)/2), so no branch discontinuities
arise in any identity proved below.

\hypertarget{exact-gamma-closed-forms}{%
\section{3. Exact Gamma Closed Forms}\label{exact-gamma-closed-forms}}

Applying Theorem 2.1 and Lemma 2.2 to each constant yields exact closed
forms. The \(\sqrt{\pi}\) factor drops from all arguments since it is
real positive.

\textbf{3.1.} c\textsubscript{1} {[}p=1, r=k{]}: factor 1+i/k gives w =
i directly.

\[c_1 = \tan\!\left(\arg\!\left(\frac{\Gamma(1+i/2)}{\Gamma(1/2+i/2)}\right)\right) \approx 0.55499611157361950342640\ldots\]

\textbf{3.2.} c\textsubscript{3} {[}p=k(k+1), r=k{]}: factor i(k+(1−i)),
i-prefactor vanishes at even N.

\[c_3 = \left|\tan\!\left(\arg\!\left(\frac{\Gamma(1-i/2)}{\Gamma(3/2-i/2)}\right)\right)\right| \approx 0.28617684964886955539522\ldots\]

\textbf{3.3.} c\textsubscript{4} {[}p=2k−1, r=k{]}: factor
(1+2i)(1−(2+i)/(5k)), even-N removes (1+2i).

\[c_4 = \left|\tan\!\left(\arg\!\left(\frac{\Gamma(4/5-i/10)}{\Gamma(3/10-i/10)}\right)\right)\right| \approx 0.24758221870861069889236\ldots\]

\textbf{3.4.} c\textsubscript{5} {[}p=k(k+1), r=k\textsuperscript{2}{]}:
factor (1+i)(1+(1+i)/(2k)), even-N removes (1+i).

\[c_5 = \left|\tan\!\left(\arg\!\left(\frac{\Gamma(5/4+i/4)}{\Gamma(3/4+i/4)}\right)\right)\right| \approx 0.20787957635076190854695\ldots\]

\textbf{3.5.} c\textsubscript{2} {[}p=(2k−1)\textsuperscript{2}, r=k{]}:
ratio p/r is quadratic. Roots of 4iz\textsuperscript{2}+(1−4i)z+i=0 are
r\textsubscript{1},\textsubscript{2} = {[}−1+4i \(\pm\)
\(\sqrt{1−8i}\){]}/(8i).

\[c_2 = \left|\tan\!\left(\arg\!\left(\sqrt{\pi} \cdot \frac{\Gamma(1-r_1/2)\cdot\Gamma(1-r_2/2)}{\Gamma(1/2-r_1/2)\cdot\Gamma(1/2-r_2/2)}\right)\right)\right| \approx 0.74328645390177619968660\ldots\]

\begin{longtable}[]{@{}
  >{\raggedright\arraybackslash}p{(\columnwidth - 8\tabcolsep) * \real{0.2000}}
  >{\raggedright\arraybackslash}p{(\columnwidth - 8\tabcolsep) * \real{0.2000}}
  >{\raggedright\arraybackslash}p{(\columnwidth - 8\tabcolsep) * \real{0.2000}}
  >{\raggedright\arraybackslash}p{(\columnwidth - 8\tabcolsep) * \real{0.2000}}
  >{\raggedright\arraybackslash}p{(\columnwidth - 8\tabcolsep) * \real{0.2000}}@{}}
\toprule\noalign{}
\begin{minipage}[b]{\linewidth}\raggedright
c
\end{minipage} & \begin{minipage}[b]{\linewidth}\raggedright
p(k), r(k)
\end{minipage} & \begin{minipage}[b]{\linewidth}\raggedright
w
\end{minipage} & \begin{minipage}[b]{\linewidth}\raggedright
Closed form (arg of)
\end{minipage} & \begin{minipage}[b]{\linewidth}\raggedright
Value
\end{minipage} \\
\midrule\noalign{}
\endhead
\bottomrule\noalign{}
\endlastfoot
c\textsubscript{1} & 1, k & i & \(\Gamma\)(1+i/2)/\(\Gamma\)(1/2+i/2) &
0.554996\ldots{} \\
c\textsubscript{2} & (2k−1)\textsuperscript{2}, k & quadratic &
\(\sqrt{\pi}\) \(\cdot\) \(\Gamma\)-product & 0.743286\ldots{} \\
c\textsubscript{3} & k(k+1), k & −i+shift &
\(\Gamma\)(1−i/2)/\(\Gamma\)(3/2−i/2) & 0.286176\ldots{} \\
c\textsubscript{4} & 2k−1, k & −(2+i)/5 &
\(\Gamma\)(4/5−i/10)/\(\Gamma\)(3/10−i/10) & 0.247582\ldots{} \\
c\textsubscript{5} & k(k+1), k\textsuperscript{2} & (1+i)/2 &
\(\Gamma\)(5/4+i/4)/\(\Gamma\)(3/4+i/4) & 0.207879\ldots{} \\
\end{longtable}

\textbf{Table 1.} The five constants and their Gamma closed forms.

\hypertarget{the-identity-s-s-4}{%
\section{\texorpdfstring{4. The Identity S\textsubscript{1} +
S\textsubscript{3} =
\(\pi\)/4}{4. The Identity S + S = /4}}\label{the-identity-s-s-4}}

\textbf{Numerical verification.} At 100-digit precision:

\[S_1 = 0.50667090321662300312194805902346766853766696879282571338879\ldots\]

\[S_3 = 0.27812742718067932321712387231630884873800968977474786328079\ldots\]

\[S_1 + S_3 = 0.78539816339744832353627193075306651628567665856757357667\ldots = \pi/4. \quad \checkmark\]

\textbf{Theorem 4.1.} S\textsubscript{1} + S\textsubscript{3} =
\(\pi\)/4, where S\textsubscript{1} =
arg(\(\Gamma\)(1+i/2)/\(\Gamma\)(1/2+i/2)) and S\textsubscript{3} =
arg(\(\Gamma\)(1−i/2)/\(\Gamma\)(3/2−i/2)).

\emph{Proof.} For the numerator product:

\[\Gamma(1+i/2)\cdot\Gamma(1-i/2) = (1/4)\cdot\Gamma(i/2)\cdot\Gamma(-i/2) = \frac{\pi}{2\sinh(\pi/2)}\]

which is real and positive (reflection formula at z = i/2). For the
denominator, \(\Gamma\)(3/2−i/2) = (1/2−i/2)\(\cdot \Gamma\)(1/2−i/2)
and the reflection formula at z = 1/2+i/2 gives
\(\Gamma(1/2+i/2)\cdot \Gamma\)(1/2−i/2) = \(\pi\)/cosh(\(\pi\)/2), so
the denominator product is (1/2−i/2)\(\cdot \pi\)/cosh(\(\pi\)/2). The
full ratio is:

\[\frac{\pi/(2\sinh(\pi/2))}{(1/2-i/2)\cdot\pi/\cosh(\pi/2)} = \coth(\pi/2)\cdot\frac{1+i}{2}\]

Since coth(\(\pi\)/2) \textgreater{} 0, the argument is arg(1+i) =
\(\pi\)/4. \(\square\)

\textbf{Corollary 4.2.} arctan(c\textsubscript{1}) +
arctan(c\textsubscript{3}) = \(\pi\)/4, giving c\textsubscript{3} =
(1−c\textsubscript{1})/(1+c\textsubscript{1}) \(\approx\)
0.28618\ldots{}

\textbf{Remark 4.3.} The algebraic relation between c\textsubscript{1}
and c\textsubscript{3} reflects the complex-conjugation symmetry of the
construction: the step polynomial ratio for c\textsubscript{3} is j+1 (a
shift of 1/j for c\textsubscript{1}), and the parameter w changes from i
to its conjugate (with a half-integer shift). The \(\pi\)/4 identity is
the angle-space manifestation of this conjugation symmetry.

\textbf{Remark 4.4.} The identity S\textsubscript{1} +
S\textsubscript{3} = \(\pi\)/4 is a special case of a general closed
form established in companion Paper 3 {[}4, Theorem 3.4{]}: for any w =
\(\pi\)/N(\(\pi\)) with \(\pi\) a Gaussian prime of norm \(\geq\) 2 (and
for the unit w = i, treated here as the degenerate N = 1 case), the
combination

\[\arg\!\left[\frac{\Gamma(1+w/2)\cdot\Gamma(1+\bar{w}/2)}{\Gamma(1/2+w/2)\cdot\Gamma(3/2+\bar{w}/2)}\right] = \arctan\!\left(\frac{b}{N+a}\right)\]

is an elementary function of (a, b, N). For w = i we have (a, b, N) =
(0, 1, 1) and arctan(1/(1+0)) = \(\pi\)/4, recovering Theorem 4.1. The
proof in {[}4{]} uses only \(\Gamma\)(z̄) = \(\Gamma\)(z)‾ and
\(\Gamma\)(3/2+z) = (1/2+z)\(\Gamma\)(1/2+z), and is algebraically
lighter than the hyperbolic-function proof presented here, though the
latter makes the geometric origin of the \(\pi\)/4 transparent.

\hypertarget{spectral-interpretation}{%
\section{5. Spectral Interpretation}\label{spectral-interpretation}}

Let D = −id/dx be the momentum operator on the circle
S\textsuperscript{1}. The zeta-regularised determinant det(D+a) =
exp(−\(\zeta\)'(0,a)) satisfies det(D+a) \(\propto\) \(\Gamma(a)^{-1}\).

\textbf{Theorem 5.1.} S\textsubscript{1} =
−arg(det(D+1/2+i/2)/det(D+1+i/2)) = Im{[}log \(\Gamma\)(1+i/2) − log
\(\Gamma\)(1/2+i/2){]}.

An equivalent form: S\textsubscript{1} = Im{[}\(\zeta\)`(0, 1+i/2) −
\(\zeta\)'(0, 1/2+i/2){]} via the Lerch formula \(\zeta\)'(0,a) = log
\(\Gamma\)(a) − ½ log(2\(\pi\)). This connects to the arithmetic
questions in §7.

\hypertarget{negative-pslq-searches}{%
\section{6. Negative PSLQ Searches}\label{negative-pslq-searches}}

Extensive PSLQ searches at the stated precision and coefficient bounds
found no relations in any of the following classes (mpmath
Bailey--Broadhurst pslq, mp.dps = 120, maxcoeff = 10\textsuperscript{6};
see Appendix A for full details):

\textbf{6.1.} Standard bases \{S\textsubscript{1}, \(\pi\), log 2,
Catalan, \(\gamma\), \(\Gamma\)(1/4), \(\Gamma\)(3/4), AGM(1,
1/\(\sqrt{2}\))\}: no relation. S\textsubscript{1}/\(\pi\),
S\textsubscript{4}/\(\pi\), S\textsubscript{5}/\(\pi\) fail algebraicity
to degree 20.

\textbf{6.2.} Im(log \(\Gamma\)(a/N+ib/M)) for N \(\leq\) 9, M \(\in\)
\{2,4,5,6,8,10\}: 273 values, no relation.

\textbf{6.3.} Im(Li\textsubscript{n}(z)) for n = 2--6 and algebraic z:
no relation.

\textbf{6.4.} Bloch-Wigner D\textsubscript{2}(z) and Zagier's
D\textsubscript{3}(z): no relation.

\textbf{6.5.} Complex Mahler measures of 1+x+y+\(\alpha\)xy: no
relation.

\textbf{6.6.} Dirichlet L-values L(1,\(\chi\)) and L(2,\(\chi\)) for
conductor \(\leq\) 12: no relation.

\textbf{6.7.} Pairwise ratios S\textsubscript{4}/S\textsubscript{1},
S\textsubscript{5}/S\textsubscript{1},
S\textsubscript{5}/S\textsubscript{4} not algebraic to degree 15.

These null results are evidence, not proof, of non-existence of
low-height algebraic relations; see Appendix A for the height-bound
context.

\hypertarget{open-questions}{%
\section{7. Open Questions}\label{open-questions}}

\textbf{(1)} Shimura/Deligne framework. Is
\(\Gamma\)(1+i/2)/\(\Gamma\)(1/2+i/2) a period ratio of a CM abelian
variety associated to \(\mathbb{Q}\)(i)?

\textbf{(2)} Generalised identity. Does S\textsubscript{4} +
S\textsubscript{5} satisfy an analogue of Theorem 4.1?

\textbf{(3)} Hurwitz zeta derivatives at complex arguments. Does
\(\zeta\)'(0,a) for complex a have arithmetic significance analogous to
the rational case?

\textbf{(4)} Infinite family. Any rational p(k)/r(k) produces a constant
\textbar tan(arg(\(\Gamma\)(1+w/2)/\(\Gamma\)(1/2+w/2)))\textbar{} for
some w \(\in\) \(\mathbb{Q}\)(i). Classifying this family is a natural
programme --- see companion Paper 3 {[}4{]}.

\hypertarget{conclusion}{%
\section{8. Conclusion}\label{conclusion}}

All five constants have exact closed forms c =
\textbar Im(z)/Re(z)\textbar{} where z =
\(\Gamma\)(1+w/2)/\(\Gamma\)(1/2+w/2). Two satisfy S\textsubscript{1} +
S\textsubscript{3} = \(\pi\)/4 (proved). S\textsubscript{1} has a
spectral interpretation as the phase of a zeta-regularised determinant.
Extensive PSLQ searches at 120-digit precision across seven classes of
known constants found no algebraic relations. Each constant is verified
absent from the OEIS, Wolfram Alpha, and ISC databases at the precision
tested.

\textbf{Python/mpmath verification} (all five closed forms, 50-digit
precision):

\begin{Shaded}
\begin{Highlighting}[]
\ImportTok{import}\NormalTok{ mpmath}\OperatorTok{;}\NormalTok{ mpmath.mp.dps }\OperatorTok{=} \DecValTok{50}
\KeywordTok{def}\NormalTok{ c(w): z }\OperatorTok{=}\NormalTok{ mpmath.gamma(}\DecValTok{1}\OperatorTok{+}\NormalTok{w}\OperatorTok{/}\DecValTok{2}\NormalTok{)}\OperatorTok{/}\NormalTok{mpmath.gamma(}\FloatTok{0.5}\OperatorTok{+}\NormalTok{w}\OperatorTok{/}\DecValTok{2}\NormalTok{)}
\ControlFlowTok{return} \BuiltInTok{abs}\NormalTok{(mpmath.im(z)}\OperatorTok{/}\NormalTok{mpmath.re(z)) }\CommentTok{\# branch{-}free form}
\NormalTok{c1 }\OperatorTok{=}\NormalTok{ c(mpmath.mpc(}\DecValTok{0}\NormalTok{,}\DecValTok{1}\NormalTok{))}\OperatorTok{;}\NormalTok{ c5 }\OperatorTok{=}\NormalTok{ c(mpmath.mpc(}\FloatTok{0.5}\NormalTok{,}\FloatTok{0.5}\NormalTok{))}
\NormalTok{S1 }\OperatorTok{=}\NormalTok{ mpmath.arg(mpmath.gamma(}\DecValTok{1}\OperatorTok{+}\OtherTok{1j}\OperatorTok{/}\DecValTok{2}\NormalTok{)}\OperatorTok{/}\NormalTok{mpmath.gamma(}\FloatTok{0.5}\OperatorTok{+}\OtherTok{1j}\OperatorTok{/}\DecValTok{2}\NormalTok{))}
\NormalTok{S3 }\OperatorTok{=}\NormalTok{ mpmath.arg(mpmath.gamma(}\DecValTok{1}\OperatorTok{{-}}\OtherTok{1j}\OperatorTok{/}\DecValTok{2}\NormalTok{)}\OperatorTok{/}\NormalTok{mpmath.gamma(}\FloatTok{1.5}\OperatorTok{{-}}\OtherTok{1j}\OperatorTok{/}\DecValTok{2}\NormalTok{))}
\BuiltInTok{print}\NormalTok{(S1 }\OperatorTok{+}\NormalTok{ S3) }\CommentTok{\# prints pi/4 to full precision}
\end{Highlighting}
\end{Shaded}

\textbf{Remark on Clifford generalisation.} The collapse mechanism
identified above --- anticommutation forces a trace-zero matrix form
whose projective action is involutive (§2.1), producing alternating
arctan sums --- is specific to Cl(2). Whether an analogous collapse
exists for higher Clifford algebras Cl(p, q) acting on spinor spaces
depends on identifying an observable whose dynamics reduce to a
tractable scalar form. We do not address this question here.

\hypertarget{appendix-a-pslq-computational-details}{%
\section{Appendix A: PSLQ Computational
Details}\label{appendix-a-pslq-computational-details}}

All PSLQ searches used the mpmath Bailey--Broadhurst implementation
(mpmath v1.3+, mp.dps = 120 with 20-digit guard, maxcoeff =
10\textsuperscript{6}). The basis sets in §6.1--6.7 were chosen as
arithmetically natural: closed under Galois conjugation over
\(\mathbb{Q}\) and spanning the known Class I--II gamma values at
denominators \(\leq\) 12. A null result at maxcoeff
10\textsuperscript{6} rules out integer relations with coefficient
height below 10\textsuperscript{6}; the Ferguson--Forcade theory {[}W.
Ferguson, R. Forcade, \emph{Math. Comp.} 40 (1983){]} implies the actual
relation height, if one exists, exceeds 10\textsuperscript{6}. Full
scripts, target values at 120-digit precision, and configuration files
are available at github.com/Loublain/cmf-gaussian-primes {[}archived:
doi:10.5281/zenodo.19670531{]}.

\hypertarget{references}{%
\section{References}\label{references}}

{[}1{]} L. Blaine. An Investigation of an Alternate Base for a
Conservative Matrix Field\ldots{} Preprint (2026).

{[}2{]} T. Raz et al.~From Euler to AI. NeurIPS 2025; arXiv:2502.17533.

{[}3{]} O. David. The conservative matrix field. arXiv:2303.09318
(2023).

{[}4{]} L. Blaine. A Family of Mathematical Constants Indexed by
Gaussian Primes. Preprint (2026).

{[}5{]} F. W. J. Olver et al.~(eds.) NIST Handbook of Mathematical
Functions. Cambridge University Press (2010). https://dlmf.nist.gov/

{[}6{]} G. Shimura. On the periods of modular forms. \emph{Mathematische
Annalen} 229, 211--221 (1977). doi:10.1007/BF01420560

{[}7{]} E. T. Whittaker, G. N. Watson. A Course of Modern Analysis. 4th
ed.~Cambridge University Press (1927).

{[}8{]} G. Raayoni et al.~Ramanujan Machine. \emph{Nature} 590, 67--73
(2021). doi:10.1038/s41586-021-03229-4

\textbf{Acknowledgements.} The author thanks Claude Sonnet 4.6
(Anthropic) for assistance with derivation, computation, and manuscript
preparation. Correspondence: louisblaine45@gmail.com

\hypertarget{footnotes}{%
\section{Footnotes}\label{footnotes}}

\textsuperscript{1} The Pauli anticommutation relations and the
anticommutative CMF condition S + T = 0 are formalized in
\texttt{SuperCMF.lean} and \texttt{PauliCMF.lean} at
github.com/Loublain/cmf-gaussian-primes. The algebraic sign flip
mechanism is proved as \texttt{AntiCMF.sign\_flip}.

\end{document}
